%!TEX root = ../main.tex
\section{MSTR vs.\ BTC ETF: Premium Decomposition}\label{sec:etf-comparison}

The launch of spot Bitcoin ETFs (e.g.\ IBIT) created a pure delta-one
BTC exposure trading at or near NAV.  Why, then, does MSTR equity
persistently trade at a premium to its own NAV?  This section answers
the question by decomposing the premium into economically distinct
components and deriving conditions under which each vehicle dominates.

%% ------------------------------------------------------------------
\subsection{The ETF as a Pure Delta-One Instrument}

\begin{definition}[BTC ETF]\label{def:etf}
A spot Bitcoin ETF with holdings $H_t^{\mathrm{ETF}}$ and shares
$N_t^{\mathrm{ETF}}$ satisfies
$$
  E_t^{\mathrm{ETF}}
  = A_t^{\mathrm{ETF}}
  = H_t^{\mathrm{ETF}}\,S_t,
  \qquad
  \pi_t^{\mathrm{ETF}} = 0.
$$
The ETF carries no debt ($D^{\mathrm{ETF}} = 0$), no leverage
($\mathrm{ILE}^{\mathrm{ETF}} = 1$), and no accumulation option
($V_{\mathrm{acc}}^{\mathrm{ETF}} = 0$).
Its equity elasticity to BTC is exactly one:
$\Delta_{\mathrm{tot}}^{\mathrm{ETF}} = 1$.
\end{definition}

\begin{remark}
In practice, ETFs trade at small, mean-reverting deviations from NAV due
to creation/redemption frictions and management fees.  We abstract from
these by setting $\pi^{\mathrm{ETF}} = 0$, which is the arbitrage-free
limit of the creation/redemption mechanism.
\end{remark}

%% ------------------------------------------------------------------
\subsection{Premium Decomposition}

\begin{theorem}[Three-Component Premium Decomposition]\label{thm:decomposition}
The MSTR log-premium admits the decomposition
$$
  \pi_t
  = \underbrace{\pi_{\mathrm{lev}}(t)}_{\text{leverage premium}}
  + \underbrace{\pi_{\mathrm{acc}}(t)}_{\text{accumulation premium}}
  + \underbrace{\pi_{\mathrm{noise}}(t)}_{\text{residual}},
$$
where
\begin{align}
  \pi_{\mathrm{lev}}(t)
  &= \log\!\bigl(\mathrm{ILE}_t\bigr)
  = \log\!\left(\frac{A_t}{A_t - D_t}\right), \label{eq:pi-lev} \\[6pt]
  \pi_{\mathrm{acc}}(t)
  &= \log\!\left(1 + \frac{V_{\mathrm{acc}}(t)}{\mathrm{NAV}_t}\right), \label{eq:pi-acc} \\[6pt]
  \pi_{\mathrm{noise}}(t)
  &= \pi_t - \pi_{\mathrm{lev}}(t) - \pi_{\mathrm{acc}}(t). \label{eq:pi-noise}
\end{align}
\end{theorem}

\begin{proof}
By definition, $E_t = e^{\pi_t}\,\mathrm{NAV}_t$.  Write equity as
the sum of leveraged NAV and the accumulation option value:
\begin{align*}
  E_t
  &= \mathrm{NAV}_t + V_{\mathrm{acc}}(t) + \varepsilon_t \\
  &= \mathrm{NAV}_t \left(1 + \frac{V_{\mathrm{acc}}(t)}{\mathrm{NAV}_t}
     + \frac{\varepsilon_t}{\mathrm{NAV}_t}\right),
\end{align*}
where $\varepsilon_t$ captures sentiment, noise, and any mispricing.
Taking logs:
$$
  \pi_t = \log\!\left(\frac{E_t}{\mathrm{NAV}_t}\right)
  = \log\!\left(1 + \frac{V_{\mathrm{acc}}(t)}{\mathrm{NAV}_t}
    + \frac{\varepsilon_t}{\mathrm{NAV}_t}\right).
$$

Now observe that the leverage premium arises because the leveraged
equity claim on BTC has higher elasticity than a dollar-for-dollar
position.  A risk-neutral investor would value the leveraged claim at
$$
  E_t^{\mathrm{lev}} = \frac{A_t}{A_t - D_t} \cdot \mathrm{NAV}_t
  = A_t,
$$
which is trivially the asset value.  However, in log-premium space,
the premium attributable purely to leverage amplification is
$$
  \pi_{\mathrm{lev}}(t) = \log(\mathrm{ILE}_t)
  = \log\!\left(\frac{A_t}{A_t - D_t}\right).
$$
This represents the excess log-return sensitivity that leverage provides
over a delta-one instrument.

The accumulation component $\pi_{\mathrm{acc}}$ captures the option
value beyond leverage, and $\pi_{\mathrm{noise}}$ is the residual.
The decomposition is exhaustive by construction.
\end{proof}

\paragraph{Calibrated decomposition.}
With $A_0 = H_0 S_0 = \$56.25$B and total senior claims
$D_0 = \$18.45$B (debt \$8.22B plus preferred liquidation value \$10.23B),
giving $\mathrm{ILE}_0 = 56.25/37.80 = 1.488$, and $\pi_0 = 0.242$:
\begin{align*}
  \pi_{\mathrm{lev}}
  &= \log(1.488) = 0.398, \\
  \pi_0 - \pi_{\mathrm{lev}}
  &= 0.242 - 0.398 = -0.156.
\end{align*}
The negative residual $\pi_0 - \pi_{\mathrm{lev}} < 0$ implies that the
current premium does not even fully compensate for leverage risk---the
accumulation option is being priced at a \emph{discount}:
$$
  \pi_{\mathrm{acc}} + \pi_{\mathrm{noise}} = -0.156.
$$
This is consistent with $\mathrm{PMRI}_0 = -0.71$ (premium below
long-run mean).

%% ------------------------------------------------------------------
\subsection{When MSTR Dominates the ETF}

\begin{proposition}[MSTR Dominance Conditions]\label{prop:mstr-dominates}
MSTR equity offers superior risk-adjusted BTC exposure relative to a
spot ETF when the following conditions hold simultaneously:
\begin{enumerate}
  \item \textbf{Positive accumulation value:} $\pi_{\mathrm{acc}}(t) > 0$,
    i.e.\ the accumulation option is in the money.
  \item \textbf{Underpriced premium:}
    $\pi_t < \pi^* \;\Leftrightarrow\; \mathrm{PMRI}_t < 0$,
    i.e.\ the current premium is below its fair (long-run) value.
  \item \textbf{Sustainable leverage:}
    $\mathrm{ILE}_t < \bar\beta$ for some prudent leverage bound
    $\bar\beta$ (e.g.\ $\bar\beta = 2$).
\end{enumerate}
\end{proposition}

\begin{proof}
We compare the expected log-return of MSTR equity to that of the ETF over
a short horizon $\Delta t$.

\textbf{ETF return.}  By definition, the ETF return equals the BTC return:
$$
  r_{\mathrm{ETF}} = \Delta\log S_t.
$$

\textbf{MSTR return.}  From the structural model:
\begin{align*}
  r_{\mathrm{MSTR}}
  &= \Delta\log E_t
  = \Delta\pi_t + \Delta\log(A_t - D_t)^+ \\
  &\approx \Delta\pi_t + \mathrm{ILE}_t \cdot \Delta\log S_t.
\end{align*}

The expected excess return of MSTR over the ETF is therefore:
$$
  \mathbb{E}[r_{\mathrm{MSTR}} - r_{\mathrm{ETF}}]
  \approx \underbrace{\mathbb{E}[\Delta\pi_t]}_{\text{premium drift}}
  + \underbrace{(\mathrm{ILE}_t - 1)}_{\text{leverage excess}}
  \cdot \mathbb{E}[\Delta\log S_t].
$$

\begin{enumerate}
  \item When $\pi_{\mathrm{acc}} > 0$, the accumulation option contributes
    positive value to equity, which is absent from the ETF.

  \item When $\mathrm{PMRI} < 0$ (premium below $\theta$), the OU dynamics
    imply $\mathbb{E}[\Delta\pi_t] = \kappa(\theta - \pi_t)\Delta t > 0$:
    the premium is expected to expand, providing a positive carry relative
    to the ETF.

  \item Leverage amplifies BTC returns ($\mathrm{ILE} > 1$) but also
    amplifies losses and increases drawdown risk.  The leverage benefit is
    sustainable only when $\mathrm{ILE}_t$ remains below a threshold
    $\bar\beta$ where bankruptcy risk is negligible (survival probability
    close to 1).
\end{enumerate}

Under all three conditions, the expected excess return is positive and
the leverage risk is bounded, making MSTR a superior vehicle.
\end{proof}

\paragraph{Current assessment.}
Applying the calibrated values:
\begin{itemize}
  \item $\mathrm{IBGR}_0 = 17.6\% > 0$: accumulation option is in the money.
  \item $\mathrm{PMRI}_0 = -0.71 < 0$: premium is below long-run mean;
    expected to expand.
  \item $\mathrm{ILE}_0 = 1.49 < 2$: leverage is moderate.
  \item Survival probability at 3 years: 91.9\%.
\end{itemize}
All three dominance conditions are satisfied.  At the current calibration
point, MSTR offers a more attractive risk-adjusted BTC exposure than the
ETF.

%% ------------------------------------------------------------------
\subsection{When the ETF Dominates MSTR}

\begin{proposition}[ETF Dominance Conditions]\label{prop:etf-dominates}
The spot BTC ETF is the superior vehicle when any of the following hold:
\begin{enumerate}
  \item \textbf{Overpriced premium:} $\pi_t > \pi^*$
    $(\mathrm{PMRI}_t > 0)$---the premium is above fair value and expected
    to compress, creating negative carry.
  \item \textbf{Excessive leverage:} $\mathrm{ILE}_t > \bar\beta$---the
    leverage amplification introduces unacceptable tail risk (survival
    probability materially below 1).
  \item \textbf{Exhausted accumulation:} $\mathrm{IBGR}_t \le 0$---the
    flywheel has stalled or reversed (e.g.\ forced BTC sales), and
    the accumulation option is out of the money.
  \item \textbf{Dividend/preferred burden:} The preferred dividend
    coverage ratio falls below a sustainability threshold, signalling
    that the capital structure is consuming value rather than creating it.
\end{enumerate}
\end{proposition}

\begin{proof}
Each condition negates a corresponding MSTR dominance condition:
\begin{enumerate}
  \item If $\mathrm{PMRI} > 0$, then $\mathbb{E}[\Delta\pi_t] < 0$:
    the premium is expected to decay, imposing a drag on MSTR returns
    that the ETF avoids entirely.

  \item If $\mathrm{ILE}_t > \bar\beta$, the leveraged claim's expected
    drawdown conditional on a BTC decline of magnitude $x$ is
    $\mathrm{ILE}_t \cdot x$, which can exceed the full BTC drop.
    The ETF's maximum drawdown is bounded by $x$.  For sufficiently
    large leverage, the Sharpe ratio of MSTR falls below that of the ETF
    despite higher expected returns.

  \item If $\mathrm{IBGR} \le 0$, the accumulation option is worthless
    ($V_{\mathrm{acc}} \le 0$), and any remaining premium represents
    pure overvaluation (or noise).

  \item Preferred dividends of \$977M/year represent a fixed claim on
    MSTR's resources.  When coverage deteriorates, the capital structure
    becomes fragile, and the risk premium demanded by common equity
    holders should rise, compressing the premium.
\end{enumerate}
In each case, the ETF provides equivalent BTC exposure without the
embedded costs and risks.
\end{proof}

%% ------------------------------------------------------------------
\subsection{The Premium as a Regime Indicator}

\begin{corollary}[Premium Regime Map]\label{cor:regime-map}
The MSTR premium regime can be characterised by the sign of PMRI and IBGR:
\begin{center}
\begin{tabular}{@{}lcc@{}}
\toprule
& $\mathrm{IBGR} > 0$ & $\mathrm{IBGR} \le 0$ \\
\midrule
$\mathrm{PMRI} < 0$ (underpriced)
  & \textbf{MSTR dominates} & Transition / distress \\
$\mathrm{PMRI} > 0$ (overpriced)
  & Hold / hedge & \textbf{ETF dominates} \\
\bottomrule
\end{tabular}
\end{center}
\end{corollary}

\paragraph{Interpretation.}
\begin{itemize}
  \item \textbf{MSTR dominates} (lower-left): The accumulation option is
    in the money and underpriced.  This is the current regime
    ($\mathrm{IBGR} = 17.6\%$, $\mathrm{PMRI} = -0.71$).

  \item \textbf{ETF dominates} (upper-right): The flywheel has stalled
    and the premium is overvalued.  Pure BTC exposure via the ETF is
    strictly preferable.

  \item \textbf{Hold/hedge} (upper-left): The accumulation is positive
    but already priced in.  The option is valuable but expensive; a
    paired position (long MSTR, short premium via options) may be
    warranted.

  \item \textbf{Transition/distress} (lower-right): The premium is
    cheap but the flywheel is broken.  This signals either a buying
    opportunity (if the flywheel restarts) or a value trap (if it
    does not).
\end{itemize}

%% ------------------------------------------------------------------
\subsection{Quantitative Comparison at Current Calibration}

We summarise the MSTR vs.\ ETF comparison using calibrated values.

\begin{table}[ht]
\centering
\caption{MSTR vs.\ spot BTC ETF: structural comparison.}
\label{tab:mstr-vs-etf}
\begin{tabular}{@{}lcc@{}}
\toprule
\textbf{Metric} & \textbf{MSTR} & \textbf{BTC ETF} \\
\midrule
BTC Delta (ILE) & 1.49 & 1.00 \\
Total Delta (TEE) & $-2.11$ & 1.00 \\
Premium ($\pi$) & 0.24 (27\% over NAV) & $\approx 0$ \\
Accumulation option & In the money & None \\
IBGR (annualised) & 17.6\% & 0\% \\
Premium carry (PMRI) & $-0.71$ (positive carry) & N/A \\
Leverage risk & Moderate (ILE $< 2$) & None \\
3-year survival & 91.9\% & 100\% \\
Preferred dividend drag & \$977M/year & None \\
Funding requirement (3y) & \$35.3B mean & None \\
\bottomrule
\end{tabular}
\end{table}

\paragraph{Bottom line.}
At the current calibration point (March 2026), MSTR offers leveraged BTC
exposure with a positive-carry accumulation option priced below its
long-run fair value.  The ETF provides simpler, lower-risk exposure
without the option component.  The indicators framework provides a
real-time dashboard for monitoring when the relative attractiveness shifts:
the key signals are PMRI crossing zero (premium becoming overpriced) and
IBGR approaching zero (flywheel stalling).
